\section{Phát biểu bài toán và phạm vi}
\label{sec:problem}

\subsection{Ký hiệu và biểu diễn quy tắc pháp luật}

Gọi \(\mathcal{R}=\{r_1,r_2,\ldots,r_n\}\) là tập các quy tắc pháp luật đã
được chuẩn hóa từ corpus. Mỗi quy tắc \(r_i\) được biểu diễn bởi bộ:

\begin{equation}
r_i =
\left(
    c_i,\,
    e_i,\,
    s_i,\,
    a_i
\right),
\label{eq:rule}
\end{equation}

trong đó:

\begin{itemize}
    \item \(c_i\) là điều kiện áp dụng của quy tắc;
    \item \(e_i\) là hệ quả pháp lý phát sinh khi điều kiện được thỏa mãn;
    \item \(s_i\) là chủ thể pháp luật liên quan;
    \item \(a_i\) là căn cứ điều luật của quy tắc.
\end{itemize}

Điều kiện \(c_i\) và hệ quả \(e_i\) được ánh xạ thành hai sự kiện pháp lý
tương ứng \(v_i^{c}\) và \(v_i^{e}\). Trong phạm vi mô hình, quy tắc \(r_i\)
định nghĩa một quan hệ có hướng:

\begin{equation}
v_i^{c}
\xrightarrow{r_i}
v_i^{e}.
\end{equation}

Quan hệ này được gọi là quan hệ điều kiện--hệ quả pháp luật. Thuật ngữ
\emph{causal} trong bài báo được sử dụng theo nghĩa thao tác hóa này: sự
kiện điều kiện kích hoạt một rule mechanism và rule mechanism làm phát sinh
sự kiện hệ quả. Quan hệ đó không được hiểu là một quan hệ nhân quả được
khám phá hoặc ước lượng từ dữ liệu quan sát.

Phiên bản corpus được sử dụng trong thực nghiệm hiện chủ yếu ánh xạ mỗi quy
tắc thành một điều kiện dương và một hệ quả dương. Các cấu trúc phức tạp hơn
như nhiều điều kiện đồng thời, ngoại lệ, phủ định, thứ tự ưu tiên giữa các
quy tắc hoặc xung đột chuẩn tắc chưa được biểu diễn đầy đủ. Giới hạn này
được thảo luận chi tiết tại Phần~\ref{sec:limitations}.

\subsection{Đồ thị nhân quả pháp luật}

Từ tập quy tắc \(\mathcal{R}\), hệ thống xây dựng một đồ thị dị thể:

\begin{equation}
\mathcal{G}
=
\left(
    \mathcal{V},
    \mathcal{E}
\right),
\end{equation}

với tập node:

\begin{equation}
\mathcal{V}
=
\mathcal{V}_{\mathrm{event}}
\cup
\mathcal{V}_{\mathrm{rule}}
\cup
\mathcal{V}_{\mathrm{article}},
\end{equation}

trong đó:

\begin{itemize}
    \item \(\mathcal{V}_{\mathrm{event}}\) chứa các sự kiện pháp lý;
    \item \(\mathcal{V}_{\mathrm{rule}}\) chứa các quy tắc đã chuẩn hóa;
    \item \(\mathcal{V}_{\mathrm{article}}\) chứa các điều luật hoặc căn cứ
    pháp lý.
\end{itemize}

Tập edge \(\mathcal{E}\) biểu diễn các quan hệ giữa event, rule và article.
Trong đó, các cạnh \texttt{CAUSES} kết nối event điều kiện với event hệ quả.
Một cặp event có thể được hỗ trợ bởi nhiều quy tắc khác nhau; vì vậy, graph
cho phép nhiều cạnh hoặc nhiều rule mechanism giữa cùng một cặp node.

Việc lưu đồng thời event, rule và article cho phép hệ thống thực hiện ba tác
vụ liên quan nhưng không đồng nhất:

\begin{enumerate}
    \item truy hồi event phù hợp với truy vấn;
    \item khôi phục causal path giữa các event;
    \item truy vết từng hop của path về rule và điều luật hỗ trợ.
\end{enumerate}

\subsection{Causal path và suy luận nhiều bước}

Một causal path có độ dài \(h\) được định nghĩa là một chuỗi có thứ tự:

\begin{equation}
p =
\left(
    v_0,\,
    r_1,\,
    v_1,\,
    r_2,\,
    \ldots,\,
    r_h,\,
    v_h
\right),
\label{eq:path}
\end{equation}

sao cho với mỗi \(j \in \{1,\ldots,h\}\), quy tắc \(r_j\) hỗ trợ quan hệ:

\[
v_{j-1}
\xrightarrow{r_j}
v_j.
\]

Trong đó:

\begin{itemize}
    \item \(v_0\) là seed event hoặc nguyên nhân khởi đầu;
    \item \(v_h\) là outcome event;
    \item \(v_1,\ldots,v_{h-1}\) là các mediator;
    \item \(r_1,\ldots,r_h\) là các quy tắc hỗ trợ từng hop.
\end{itemize}

Trong cấu hình thực nghiệm chính, hệ thống truy hồi các path có tối đa hai
hop, tức là \(h \leq 2\). Một đường hai bước có dạng:

\[
v_0
\xrightarrow{r_1}
v_1
\xrightarrow{r_2}
v_2,
\]

trong đó \(v_1\) vừa là hệ quả của quy tắc \(r_1\), vừa là điều kiện của
quy tắc \(r_2\).

Giới hạn hai hop được lựa chọn vì benchmark hiện tại chủ yếu được xây dựng
từ các chuỗi hai bước có trong graph. Đây không phải giới hạn lý thuyết của
biểu diễn path; tuy nhiên, khi tăng số hop, không gian ứng viên và khả năng
tích lũy lỗi truy hồi cũng tăng đáng kể.

\subsection{Phân biệt quan hệ trực tiếp và quan hệ gián tiếp}

Cho hai event \(A\) và \(C\), một quan hệ trực tiếp tồn tại khi graph chứa
ít nhất một rule mechanism:

\[
A \xrightarrow{r} C.
\]

Ngược lại, nếu graph chỉ chứa một đường:

\[
A \xrightarrow{r_1} B \xrightarrow{r_2} C,
\]

thì hệ thống có thể kết luận rằng \(A\) liên quan đến \(C\) thông qua
mediator \(B\), nhưng không được kết luận rằng tồn tại một cạnh trực tiếp
\(A \rightarrow C\).

Sự phân biệt này đặc biệt quan trọng đối với các câu hỏi như:

\begin{quote}
``\(A\) có trực tiếp dẫn đến \(C\) hay không?''
\end{quote}

Nếu graph chỉ hỗ trợ đường hai bước qua \(B\), claim trực tiếp phải bị bác
bỏ, ngay cả khi \(C\) vẫn có thể được suy ra gián tiếp từ \(A\).

\subsection{Truy vấn và các loại claim}

Gọi \(q\) là một truy vấn ngôn ngữ tự nhiên. Hệ thống phân tích \(q\) để
xác định loại quan hệ mà người dùng yêu cầu. Các loại claim chính gồm:

\begin{description}
    \item[\texttt{STANDARD\_CAUSAL}] Truy vấn yêu cầu xác định hoặc giải
    thích một quan hệ điều kiện--hệ quả thông thường.

    \item[\texttt{DIRECT\_CAUSAL\_CLAIM}] Truy vấn khẳng định hoặc hỏi về
    sự tồn tại của một cạnh trực tiếp giữa hai event.

    \item[\texttt{REMOVE\_MEDIATOR\_OUTCOME\_REMAINS}] Truy vấn cho rằng
    outcome vẫn xuất hiện khi mediator bị loại bỏ.

    \item[\texttt{REMOVE\_MEDIATOR\_OUTCOME\_DISAPPEARS}] Truy vấn cho rằng
    outcome biến mất khi mediator bị loại bỏ.

    \item[\texttt{MEDIATOR\_NECESSARY\_CLAIM}] Truy vấn hỏi liệu mediator
    có cần thiết đối với outcome hay không.

    \item[\texttt{COUNTERFACTUAL\_UNSPECIFIED}] Truy vấn chứa tín hiệu phản
    thực tế nhưng chưa chỉ rõ outcome được kỳ vọng giữ nguyên hay biến mất.
\end{description}

Việc xác định claim chỉ dựa trên nội dung truy vấn và cấu trúc sự kiện được
truy hồi. Hệ thống không sử dụng \texttt{question\_type}, nhãn quyết định
chuẩn, gold path hoặc gold answer của benchmark.

\subsection{Factual world và counterfactual world}

Mỗi event \(v \in \mathcal{V}_{\mathrm{event}}\) được liên kết với một biến
trạng thái:

\begin{equation}
X_v
\in
\left\{
    \mathrm{TRUE},
    \mathrm{FALSE},
    \mathrm{UNKNOWN}
\right\}.
\end{equation}

Ba trạng thái được diễn giải như sau:

\begin{itemize}
    \item \(\mathrm{TRUE}\): event được suy ra là xảy ra trong world đang
    xét;
    \item \(\mathrm{FALSE}\): event được suy ra là không xảy ra;
    \item \(\mathrm{UNKNOWN}\): chưa đủ căn cứ để xác định hoặc có xung đột
    giữa các cơ chế suy luận.
\end{itemize}

Factual world \(W^{F}\) được xác định từ seed event và các rule mechanism
được kích hoạt trong graph. Với một causal path:

\[
A \rightarrow M \rightarrow Y,
\]

hệ thống trước hết kiểm tra các trạng thái factual:

\[
X_A^{F},\quad X_M^{F},\quad X_Y^{F}.
\]

Để đánh giá vai trò của mediator \(M\), LegalSCM thực hiện hard
intervention:

\begin{equation}
do
\left(
    X_M = \mathrm{FALSE}
\right).
\label{eq:do}
\end{equation}

Theo semantics của can thiệp trong SCM
\cite{pearl2009causality}, phép toán trong
Phương trình~\ref{eq:do} không chỉ thay đổi giá trị của \(M\). Nó còn thay
thế cơ chế cấu trúc thông thường tạo ra \(M\), tương đương với việc cắt các
cơ chế đầu vào của mediator trong world được can thiệp.

Counterfactual world tương ứng được ký hiệu:

\[
W^{CF}_{M}
=
W^{F}
\mid
do
\left(
    X_M = \mathrm{FALSE}
\right).
\]

Sau can thiệp, LegalSCM suy luận lại trạng thái của outcome:

\[
X_Y^{CF}
=
X_Y
\left(
    W^{CF}_{M}
\right).
\]

Việc so sánh \(X_Y^{F}\) và \(X_Y^{CF}\) cho phép phân biệt các trường hợp:

\begin{itemize}
    \item nếu
    \(X_Y^{F}=\mathrm{TRUE}\) và
    \(X_Y^{CF}=\mathrm{FALSE}\),
    mediator được xem là cần thiết trong mô hình đang xét;

    \item nếu
    \(X_Y^{F}=\mathrm{TRUE}\) và
    \(X_Y^{CF}=\mathrm{TRUE}\),
    outcome vẫn có thể được duy trì sau can thiệp;

    \item nếu một trong hai trạng thái là
    \(\mathrm{UNKNOWN}\),
    kết quả can thiệp chưa đủ để đưa ra kết luận chắc chắn.
\end{itemize}

Khái niệm ``cần thiết'' ở đây chỉ có ý nghĩa tương đối với các rule mechanism
và factual context được đưa vào LegalSCM. Nó không chứng minh tính cần thiết
pháp lý trong mọi tình huống thực tế.

\subsection{Path ablation}

Để tạo baseline, nghiên cứu sử dụng một phép path ablation. Với cùng đường:

\[
A \rightarrow M \rightarrow Y,
\]

mediator \(M\) được loại khỏi graph, sau đó hệ thống tìm các đường thay thế
từ \(A\) đến \(Y\). Nếu không còn đường phù hợp, mediator được xem là cần
thiết theo tiêu chí reachability. Nếu tồn tại đường thay thế đạt ngưỡng
điểm, outcome được xem là vẫn có thể tiếp cận.

Path ablation và hard intervention giải quyết hai bài toán khác nhau:

\begin{itemize}
    \item path ablation hỏi liệu graph sau khi xóa \(M\) còn chứa đường từ
    \(A\) đến \(Y\) hay không;

    \item LegalSCM hỏi trạng thái của \(Y\) thay đổi như thế nào khi
    \(M\) bị gán cưỡng bức về \(\mathrm{FALSE}\) và các cơ chế liên quan
    được suy luận lại.
\end{itemize}

Vì vậy, kết quả path ablation không được diễn giải như kết quả của phép
\(do(\cdot)\). Trong nghiên cứu này, path ablation chỉ được sử dụng làm
baseline để đánh giá chênh lệch về chất lượng, chi phí và mức độ chi tiết
của trace.

\subsection{Quyết định xác minh}

Với truy vấn \(q\), tập path ứng viên \(\mathcal{P}_q\) và kết quả phân tích
claim, hệ thống tạo một quyết định:

\begin{equation}
D_q
\in
\left\{
    \texttt{SUPPORTED},
    \texttt{REJECT\_DIRECT\_CLAIM},
    \texttt{UNCERTAIN}
\right\}.
\end{equation}

Các nhãn được diễn giải như sau:

\begin{description}
    \item[\texttt{SUPPORTED}] Cấu trúc graph và kết quả xác minh hỗ trợ claim
    được phát hiện trong truy vấn.

    \item[\texttt{REJECT\_DIRECT\_CLAIM}] Claim trực tiếp hoặc claim cần bác
    bỏ không phù hợp với cấu trúc path và kết quả xác minh. Ví dụ, truy vấn
    khẳng định \(A\) trực tiếp dẫn đến \(C\), nhưng graph chỉ hỗ trợ
    \(A \rightarrow B \rightarrow C\).

    \item[\texttt{UNCERTAIN}] Hệ thống không tìm được path đủ tin cậy, factual
    path không thể được xác nhận hoặc kết quả can thiệp chứa trạng thái chưa
    xác định.
\end{description}

Ngoài nhãn \(D_q\), hệ thống lưu một decision score, phần giải thích, primary
path, loại claim và các trường factual/counterfactual liên quan. Các trường
này được chuyển nguyên vẹn sang bước sinh câu trả lời nhằm hạn chế trường
hợp answer generator đưa ra kết luận trái với bước xác minh.

\subsection{Dạng câu trả lời}

Loại claim xác định semantics của quyết định, trong khi answer intent xác
định dạng thông tin cần trình bày. Hệ thống hỗ trợ các answer intent sau:

\begin{description}
    \item[\texttt{FINAL\_OUTCOME}] Trả về hệ quả cuối cùng của primary path.

    \item[\texttt{INITIAL\_CAUSE}] Trả về nguyên nhân hoặc sự kiện khởi đầu.

    \item[\texttt{BRIDGE\_EVENT}] Trả về mediator nằm giữa nguyên nhân và
    outcome.

    \item[\texttt{CHAIN\_EXPLANATION}] Trình bày lần lượt từng hop cùng căn
    cứ điều luật tương ứng.

    \item[\texttt{YES\_NO}] Trả lời có, không hoặc chưa đủ căn cứ, phù hợp
    với quyết định xác minh.

    \item[\texttt{BRANCH}] Trình bày nhiều outcome có chung một prefix nhân
    quả.

    \item[\texttt{CONVERGENCE}] Trình bày nhiều điều kiện ban đầu hội tụ vào
    cùng một suffix hoặc outcome.

    \item[\texttt{CLAIM\_VERIFICATION}] Trả lời trực tiếp về tính đúng hoặc
    sai của claim direct/counterfactual.
\end{description}

Answer intent được suy ra từ truy vấn và kết quả query analysis của bước xác
minh. Tương tự claim detection, quá trình này không đọc nhãn loại câu hỏi
hoặc gold answer.

\subsection{Đầu ra của bài toán}

Toàn bộ pipeline có thể được mô tả như ánh xạ:

\begin{equation}
F:
\left(
    q,\,
    \mathcal{R},\,
    \mathcal{G}
\right)
\longrightarrow
\left(
    \widehat{\mathcal{E}}_q,\,
    \widehat{\mathcal{P}}_q,\,
    D_q,\,
    A_q,\,
    C_q,\,
    M_q
\right),
\end{equation}

trong đó:

\begin{itemize}
    \item \(\widehat{\mathcal{E}}_q\) là tập evidence được truy hồi;
    \item \(\widehat{\mathcal{P}}_q\) là các causal path ứng viên và primary
    path được chọn;
    \item \(D_q\) là quyết định xác minh;
    \item \(A_q\) là câu trả lời cuối cùng;
    \item \(C_q\) là tập citation được sử dụng;
    \item \(M_q\) là metadata về truy hồi, factual world, intervention,
    counterfactual world và quá trình sinh câu trả lời.
\end{itemize}

Một đầu ra được xem là tốt khi đồng thời thỏa mãn nhiều tiêu chí:

\begin{enumerate}
    \item rule và event liên quan được truy hồi;
    \item primary path phản ánh đúng thứ tự suy luận;
    \item quyết định xác minh phù hợp với loại claim;
    \item câu trả lời đúng answer intent;
    \item citation thực sự hỗ trợ các hop được sử dụng;
    \item metadata cho phép truy vết và tái lập quá trình xử lý.
\end{enumerate}

\subsection{Phạm vi và các giả định}

Nghiên cứu được thực hiện dưới các giả định và giới hạn sau:

\begin{enumerate}
    \item \textbf{Quan hệ quy phạm thay cho causal discovery.}
    Hướng cạnh được lấy từ cấu trúc điều kiện--hệ quả của quy tắc pháp luật,
    không được học từ dữ liệu quan sát.

    \item \textbf{Suy luận tất định ba giá trị.}
    LegalSCM không biểu diễn xác suất. Các event được gán
    \(\mathrm{TRUE}\), \(\mathrm{FALSE}\) hoặc \(\mathrm{UNKNOWN}\).

    \item \textbf{Can thiệp giới hạn.}
    Thực nghiệm hiện tại chỉ tập trung vào việc đặt mediator về
    \(\mathrm{FALSE}\). Các can thiệp đa biến hoặc can thiệp trên giá trị
    liên tục nằm ngoài phạm vi nghiên cứu.

    \item \textbf{Đường dẫn ngắn.}
    Cấu hình đánh giá chính giới hạn causal path ở tối đa hai hop.

    \item \textbf{Biểu diễn quy tắc đơn giản hóa.}
    Corpus hiện chưa bao phủ đầy đủ ngoại lệ, phủ định, nhiều điều kiện đồng
    thời và ưu tiên giữa các quy tắc.

    \item \textbf{Sinh câu trả lời extractive.}
    Thực nghiệm chính sử dụng answer generator theo mẫu, không đánh giá khả
    năng diễn đạt của một mô hình ngôn ngữ sinh tự do.

    \item \textbf{Benchmark cùng nguồn graph.}
    Các câu hỏi được sinh từ graph dùng trong pipeline. Vì vậy, kết quả đánh
    giá phản ánh khả năng khôi phục và xử lý cấu trúc đã mã hóa, không phản
    ánh đầy đủ khả năng tổng quát hóa sang corpus pháp luật độc lập.

    \item \textbf{Không thay thế tư vấn pháp lý.}
    Hệ thống được xây dựng cho mục đích nghiên cứu và không được sử dụng như
    một nguồn tư vấn hoặc kết luận pháp lý chính thức.
\end{enumerate}
